% ---------------------------------------------------------------------------------------------
% Appendix K: what the opponent detected (main study).
% Needs: figures/fig_detection_content.pdf. Cross-references used: fig:results (main Figure 2), app:opponent.
% Built by figures/channel_activation/fig_detection_content.py; all numbers are in detection_content.csv.
% The \setcounter line forces the letter K; delete it if this is already the eleventh appendix section.
% ---------------------------------------------------------------------------------------------
\setcounter{section}{10}
\section{What the opponent detected}\label{app:detection}\setcounter{table}{0}\setcounter{figure}{0}

The main text infers the source of each model's losses from its behavior: marginal bias and runs statistics (Figure~\ref{fig:results}) and the Brown--Rosenthal regressions. Proteus's logged state gives a more direct account, because on every round it records which channels were significant, which channel it acted on, and which action it predicted (Appendix~\ref{app:opponent}). We used these records to classify each of the 19,600 main-study rounds by the content of the opponent's prediction. A round had \emph{no detection} if no channel was significant, in which case Proteus chose at random. It was a \emph{side-bias} round if the marginal channel was significant and the prediction Proteus acted on was the side that channel favored, whichever channel supplied it. Among the remaining rounds, Proteus either \emph{predicted a switch} from the model's previous choice or \emph{predicted a repeat} of it. We classify by the prediction because the channels themselves are not behaviors: the history channels respond to alternation and to repetition alike, a single alternating pattern makes all eight of them significant at once, and a strong side bias is visible to them as well as to the marginal channel. The classification is descriptive and was defined after the data were collected.

Figure~\ref{fig:detection} shows that the opponent caught two different regularities under the informed prompt, and that it profited from both. For Sonnet and Haiku it predicted a switch on 89\% and 84\% of rounds and won 82\% and 76\% of them, and the side-bias category was empty. For GPT-5.4, GPT-5.4-mini, and Gemini 3.1 Flash Lite the pattern was the reverse: 83\% to 86\% of rounds were side-bias rounds, on which Proteus won 60\% to 68\% of the time, and predicted switches were nearly absent. Opus showed both (side bias on 21\% of rounds, predicted switches on 47\%). Weighting each category by its frequency and by Proteus's margin over one-half, predicted switches account for 57\% of the seven models' combined shortfall below a .50 win rate under the informed prompt and side bias for 40\%. An opponent that detected only alternation would therefore have had almost nothing to act on against three of the six models whose win rates declined. Under the minimal prompt the opponent was idle on half of all rounds, and when it acted it almost always predicted a repeat (45\% of rounds, 60\% won); these streak-following predictions account for essentially all of the smaller minimal-prompt shortfall and were most frequent for Opus and Gemini 3 Flash Preview (76\% and 72\% of rounds). Informed prompting thus changed what the opponent found, from streaks to sustained alternation or side bias, and not merely how often it found something.

Two further points bear on interpretation. First, on rounds with no detection Proteus won 49.4\% of the time under the minimal prompt and 51.9\% under the informed prompt, so the models' losses arose on the rounds where the opponent had identified a regularity and not from an advantage built into the game. Second, these figures describe the sessions as they were played. Each model's later choices depended on the opponent's earlier ones, so the shares are an accounting of where the realized losses occurred, not an estimate of how the same models would have fared against an opponent restricted to one kind of detection.

\begin{figure}[htbp]
\centering
\includegraphics{figures/fig_detection_content.pdf}
\caption{What Proteus detected in the main study, by prompt. Each of a model's 1,400 rounds per prompt falls in one column: no channel significant (Proteus chose at random); side bias (the marginal channel was significant and Proteus predicted the side it favored); or, otherwise, a prediction that the model would switch from or repeat its previous choice. Side bias is a marginal regularity; predicted switches and repeats are sequential ones. Shading and the upper number give the percentage of rounds in the category. The lower number, in bold, is the percentage of those rounds that Proteus won, shown where the category has at least 30 rounds; independent 50/50 play holds it to 50. Models are in the order of Figure~\ref{fig:results}.}
\label{fig:detection}
\end{figure}
